% Pursuing Stacks (cotes 134-2 to 134-8): archive cote 134-2, site page 54 (PDF page 55).
% appendix (letters to Breen). Draft compared: ps1.tex:1356-1385.
% AI-transcribed by Claude Opus 5.5 (claude-opus-5-5, adaptive thinking, effort high), Message Batches API, batch ps1 (msgbatch_01Rv5hCaNTtUt4JHx1yqoHPL), 2026-11-12; not yet checked by eye.
\scanpage{134-2/54}
\mspage{\hand{39}} \hfill \mspage{7}

\[ f_n : \hand{\mathcal{C}_n \longrightarrow \mathcal{D}_n} . \]

On doit avoir alors:

(A') f$_n$est une n-équivalence sss le n- foncteur $\hand{\varphi \longrightarrow \varphi}\circ f_n$
\[ f_n^*: \quad \underline{Hom}\, (\hand{\mathcal{D}}_n, \ ((n-1)\text{-}Cat) \hand{\longrightarrow} \underline{Hom}(\hand{\mathcal{C}}_n, \ ((n-1)\text{-}Cat)) \]
allant des (n-1) systèmes locaux sur $\hand{\mathcal{D}}_n$ (ou $\hand{\mathcal{D}}_{n+1}$ , c'est pareil) vers
les (n-1)- systèmes locaux sur $\hand{\mathcal{C}}_n$, est une n-équivalence.

(B') $f_n$ est une n-équivalence sss pour tout n-système local $\hand{\mathcal{F}}$ sur
$\hand{\mathcal{D}}_{n+1}$ ,
\[ \hand{\mathcal{F}} : \hand{\mathcal{D}}_{n+1} \hand{\longrightarrow} (n\text{-}Cat), \]
le n-foncteur induit par $f_{n+1}$
\[ \underbrace{\underline{Hom}\, ( e_n^{\hand{\mathcal{D}_{n+1}}}, \ \hand{\mathcal{F}} )}_{\hand{\parallel\ \text{déf}} \atop \hand{\Gamma_{\mathcal{D}_{n+1}}(\mathcal{F})}} \hand{\longrightarrow} \underbrace{\underline{Hom}\ (e_n^{\hand{\mathcal{C}_{n+1}}} \ , \ f^*_{n+1} \hand{\mathcal{F}} \ )}_{\hand{\parallel\ \text{déf}} \atop \hand{\Gamma_{\mathcal{C}_{n+1}}(\mathcal{F})}} \]
\margin{\unsure{espacer} (circled, with a vertical double arrow)}
est une n-équivalence.

\margin{\unsure{(App.8)|(Arr.8)}} \hand{(} La construction de la cohomologie d'un topos en termes d'intégration
des champs ne fait aucun appel à la notion de complexe de faisceaux abéliens,
encore moins à la technique des résolutions injectives. On a l'impression
que dans son esprit, via la définition (qui reste à expliciter ! ) des
n-champs, elle s'apparenterait plutôt aux calculs ``Cechistes'' en termes
d'hyperrecouvrements. Or ces derniers se décrivent à l'aide d'une petite
dose d'algèbre semi-simpliciale. Je ne sais si une théorie des champs
et opérations sur iceux pourra s'écrire \uline{sans} jamais utiliser de
l'algèbre semi-simpliciale. Si oui, cela ferait essentiellement trois
approches distinctes pour construire la cohomologie d'un topos \unsure{\over{;}{:}|:}

a) point de vue des complexes de faisceaux, des résolutions injectives,
des catégories dérivées (\uline{algèbre homologique commutative})\hand{,}

b) point de vue \over{C}{\v{C}}echiste ou semi-simplicial (\uline{algèbre homotopique})\hand{,}

c) point de vue des n-champs (algèbre catégorique, ou \uline{algèbre homo-}%

\mspage{\hand{54}}
